%% ============================================================
%%  SECCIÓN 3 — Objetivos
%%  Responsable: [Nombre]
%% ============================================================

\section{Objetivos}

\subsection{Objetivo general}

Determinar experimentalmente la relación entre el esfuerzo aplicado y la deformación
unitaria en un resorte metálico y una tira de jebe, verificando la validez de la Ley de
Hooke y calculando el módulo de Young de cada material mediante análisis gráfico y
ajuste lineal.

\subsection{Objetivos específicos}

\begin{itemize}
    \item Medir la longitud natural, el diámetro de la sección transversal y la masa
    del resorte y la liga en condiciones sin carga.

    \item Registrar las deformaciones longitudinales y transversales producidas por
    seis cargas acumulativas sucesivas, tanto en la etapa de carga como en la de
    descarga.

    \item Calcular el esfuerzo técnico $\sigma_0 = F/S_0$, el esfuerzo real
    $\sigma = F/S$ y la deformación unitaria $\varepsilon = \Delta\ell/\ell_0$
    para cada configuración de carga.

    \item Construir las gráficas $F$ vs.\ $\Delta\ell$ y $\sigma$ vs.\ $\varepsilon$
    para el resorte y la liga, analizando las fases de carga y descarga.

    \item Determinar la constante recuperadora del resorte $k$ y el módulo de Young $Y$
    a partir de la pendiente del ajuste lineal en la región elástica.

    \item Estimar numéricamente el trabajo realizado para deformar el resorte desde su
    posición de equilibrio hasta la tercera carga, mediante integración numérica de la
    curva $F$ vs.\ $\Delta\ell$.

    \item Estimar el área encerrada por el lazo de histéresis de la liga en el
    diagrama $\sigma$--$\varepsilon$, cuantificando la energía disipada por
    unidad de volumen durante el ciclo de carga--descarga.

    \item Comparar el comportamiento elástico del resorte con el de la liga, analizando
    la histéresis observada en el ciclo carga--descarga.
\end{itemize}